\documentclass[11pt,a4paper]{article}

% --- Packages ---
\usepackage[margin=.75in]{geometry}
\usepackage[utf8]{inputenc}
\usepackage[T1]{fontenc}
\usepackage{lmodern} % Fix for font not found error
\usepackage{titlesec}
\usepackage{enumitem}
\usepackage{hyperref}
\usepackage{xcolor}
\usepackage{fancyhdr}
\usepackage{comment}
\usepackage{cleveref}
\usepackage[bottom]{footmisc}
\usepackage{graphicx}
\usepackage{longtable}
\usepackage{booktabs}
\usepackage{array}
\usepackage{calc}
\usepackage{etoolbox}

% Standard Pandoc setup
\providecommand{\tightlist}{%
  \setlength{\itemsep}{0pt}\setlength{\parskip}{0pt}}

\crefformat{section}{\S#2#1#3}
\crefformat{subsection}{\S#2#1#3}
\crefformat{subsubsection}{\S#2#1#3}
\crefformat{figure}{#2Figure~#1#3}
\crefformat{footnote}{#2\footnotemark[#1]#3}

\linespread{0.9} % Riduce lo spazio tra le linee

% --- Colors ---
\definecolor{primary}{RGB}{0, 0, 0}
\definecolor{links}{HTML}{0000EE}
\definecolor{blue}{HTML}{26428B}

\hypersetup{
    colorlinks=true,
    linkcolor=links,
    urlcolor=links,
    citecolor=links,
}

% --- Section Formatting ---
\titleformat{\section}
  {\normalfont\Large\bfseries\color{blue}}{\thesection}{1em}{}
\titleformat{\subsection}
{\normalfont\large\bfseries\color{blue}}{\thesubsection}{1em}{}
\titleformat{\subsubsection}
{\normalfont\bfseries\color{blue}}{\thesubsubsection}{1em}{}

% --- Header and Footer ---
\pagestyle{fancy}
\fancyhead[L]{\small Giuseppe Capasso}
\fancyhead[R]{\small intro-to-linux}
\fancyfoot[C]{\thepage}
\renewcommand{\headrulewidth}{0.4pt}

% --- Custom Title Command ---
\newcommand{\proposaltitle}[1]{
    \begin{center}
        \vspace*{0.1cm}
        {\LARGE \bfseries \color{blue} #1} \\[1.5ex]
        {\large \textbf{Giuseppe Capasso}} \\[1ex]
                \small{
            \vspace{0.1cm}
            \textbf{Materia:} intro-to-linux\\
        }
            \end{center}
    \vspace{0.3cm}
    \hrule height 0.5pt
    \vspace{0.5cm}
}

\begin{document}

\proposaltitle{Obiettivo: Mettere in pratica tutte le competenze acquisite nel corso (CLI, Storage, Permessi, Rete,}

Obiettivo: Mettere in pratica tutte le competenze acquisite nel corso
(CLI, Storage, Permessi, Rete, Servizi, Sicurezza) configurando un
server di produzione da zero. Durata stimata: 2 Ore Scenario: Sei stato
assunto come Junior System Administrator presso la startup ``TechNova''.
Ti è stata consegnata una macchina virtuale nuda (Ubuntu Server) e un
secondo disco vuoto. Devi preparare l'infrastruttura per ospitare il
nuovo sito aziendale in modo sicuro e resiliente.

Fase 1: Accesso e Sicurezza Iniziale (Lezione 2 \& 5) Il server non deve
mai essere gestito direttamente dalla console di VirtualBox, e l'accesso
deve essere blindato. 1. Configura il Port Forwarding su VirtualBox per
la porta SSH (es. 2222 -\textgreater{} 22) e per il traffico Web (es.
8080 -\textgreater{} 80). 2. Collegati via SSH dal tuo terminale
Windows/Mac locale. 3. Genera e copia le tue chiavi SSH sul server. 4.
Modifica /etc/ssh/sshd\_config per disabilitare l'accesso con password
(PasswordAuthentication no). Riavvia il servizio SSH. 5. Configura il
Firewall (UFW): • Blocca tutto in ingresso. • Consenti SSH e HTTP. •
Attiva UFW e verifica lo stato.

Fase 2: Provisioning dello Storage (Lezione 3) Il sito aziendale
produrrà molti dati. Non possiamo salvarli sul disco di sistema, ci
serve il disco secondario. 1. Aggiungi un disco da 5GB alla VM su
VirtualBox (se non lo hai già fatto). 2. Trova il nuovo disco usando
lsblk. 3. Usa cfdisk per creare un'unica partizione primaria che occupi
tutto lo spazio. 4. Formatta la partizione in ext4. 5. Crea la cartella
/var/www/technova. 6. Ottieni l'UUID del nuovo disco e configura
/etc/fstab affinché la nuova partizione venga montata in automatico su
/var/www/technova. 7. Testa il mount con sudo mount -a e verifica con df
-h.

Fase 3: Deploy dell'Applicazione (Lezione 2 \& 4) Installiamo il server
web e prepariamo i permessi. 1. Installa il server web Nginx: sudo apt
update \&\& sudo apt install nginx. 2. Cambia il proprietario della
cartella /var/www/technova: assegnala all'utente www-data e al gruppo
www-data. 3. Imposta i permessi della cartella a 755. 4. Crea un file
index.html all'interno di /var/www/technova contenente il messaggio:

Benvenuti in TechNova - Server Sicuro

. 5. Trick avanzato: Di default, Nginx legge da /var/www/html. Rinomina
la cartella di default e crea un link simbolico, oppure semplicemente
modifica la configurazione di default di Nginx
(/etc/nginx/sites-available/default) cambiando root /var/www/html; in
root /var/www/technova;. 6. Riavvia Nginx con Systemd.

Fase 4: Servizio di Sicurezza Personalizzato (Lezione 5) L'azienda vuole
un sistema che registri l'orario di accensione del server su un file di
testo ogni volta che la macchina fa il boot. 1. Crea uno script
semplicissimo in /usr/local/bin/boot\_logger.sh: Bash \#!/bin/bash echo
``Server TechNova avviato alle: \$(date)'' \textgreater\textgreater{}
/var/log/technova\_boot.log

\begin{enumerate}
\def\labelenumi{\arabic{enumi}.}
\setcounter{enumi}{1}
\tightlist
\item
  Rendi lo script eseguibile (chmod +x).
\item
  Crea una Unit Systemd chiamata bootlogger.service che esegua questo
  script al boot (usa Type=oneshot).
\item
  Abilita il servizio.
\end{enumerate}

Fase 5: Strategia di Backup (Lezione 3) Il sito deve essere salvato
periodicamente. 1. Crea la cartella /backup. 2. Usa tar per creare un
archivio compresso (.tar.gz) dell'intera cartella /var/www/technova e
salvalo dentro /backup.

3. Bonus: Inserisci il comando tar in crontab -e (cerca su internet come
si fa!) per farlo eseguire automaticamente ogni notte a mezzanotte (0 0
* * *). IL TEST FINALE: Apri il browser sul tuo PC fisico e vai su
http://127.0.0.1:8080. Se vedi la pagina di TechNova, hai completato con
successo l'infrastruttura!



\end{document}
